\documentclass{article}

\PassOptionsToPackage{numbers, sort, compress}{natbib}
\usepackage[main,final]{neurips_2025}

\usepackage{tempora}
\usepackage[T1,T2A]{fontenc}
\usepackage[utf8]{inputenc}
\usepackage[russian,english]{babel}
\usepackage{hyperref}
\usepackage{url}
\usepackage{booktabs}
\usepackage{nicefrac}
\usepackage{microtype}
\usepackage{xcolor}

%%%

\usepackage{subcaption}
\usepackage{graphicx}
\usepackage{multirow}
\usepackage{amsmath,amssymb,amsfonts}
\usepackage{amsthm}
\usepackage{mathrsfs}
\usepackage{xcolor}
\usepackage{textcomp}
\usepackage{manyfoot}
\usepackage{algorithm}
\usepackage{algorithmicx}
\usepackage{algpseudocode}
\usepackage{listings}

\newtheorem{theorem}{Theorem} % continuous numbers
%%\newtheorem{theorem}{Theorem}[section] % sectionwise numbers
%% optional argument [theorem] produces theorem numbering sequence instead of independent numbers for Proposition
\newtheorem{proposition}[theorem]{Proposition}% 
\newtheorem{lemma}{Lemma}% 
%%\newtheorem{proposition}{Proposition} % to get separate numbers for theorem and proposition etc.

\newtheorem{example}{Example}
\newtheorem{remark}{Remark}

\newtheorem{definition}{Definition}
\newtheorem{assumption}{Assumption}

%%%


\title{Title}


% The \author macro works with any number of authors. There are two commands
% used to separate the names and addresses of multiple authors: \And and \AND.
%
% Using \And between authors leaves it to LaTeX to determine where to break the
% lines. Using \AND forces a line break at that point. So, if LaTeX puts 3 of 4
% authors names on the first line, and the last on the second line, try using
% \AND instead of \And before the third author name.


\author{
  Михаил Плигин\\
  MIPT\\
  Moscow, Russia\\
  \texttt{email@phystech.edu}\\
  \And
  Name Surname\\
  MIPT\\
  Moscow, Russia\\
  \texttt{email@phystech.edu}\\
}


\begin{document}


\maketitle

\begin{abstract}
    \textit{Актуальность.} Декодеры речи на глубоких нейросетях выигрывают в точности ценой прозрачности: по их весам трудно проверить, опирается ли решение на нейронную активность, а не на мышечные и акустические артефакты, что критично для клинических интерфейсов мозг–компьютер.
    \textit{Цель работы.} Предложить архитектуру для классификации слов по внутричерепной электроэнцефалографии, которая сохраняет интерпретируемый блок детектора огибающей из работы \cite{Petrosyan2022} и превосходит исходную модель по точности. 
    \textit{Методы.} В работе рассматривается модель с обучаемыми пространственными и полосовыми фильтрами, выпрямлением и сглаживанием, веса которых переводятся в пространственные и частотные паттерны, — а двунаправленный блок долгой краткосрочной памяти заменён причинной временной свёрточной сетью с расширенными  свёртками. Модель обучалась по двухэтапной схеме исходной работы: регрессия лог-мел-спектрограммы речи, затем классификация 26 слов и тишины по промежуточным признакам. Сравнение с воспроизведённой исходной моделью и её вариантом без временного модуля проведено на данных двух пациентов, и открытых данных \cite{Verwoert2022}. 
    \textit{Результаты.} Предложенная модель достигает точности $[X1] \%$ и $[X2] \%$ (для 1 и 2 пациента соответственно) против $[B1]\%$ и $[B2]\%$ у воспроизведённой исходной архитектуры при уровне случайного угадывания $3,7\%$. Пространственные и частотные паттерны блок детектора огибающей после замены временного модуля сохраняются и указывают на [область коры] и диапазон $[f_1,f_2]Hz$.
    \textit{Значимость.} Интерпретируемость первых слоёв, не ограничивает точность: резерв качества компактного декодера речи лежит во временном модуле, и его замена повышает точность без потери возможности проверить нейронное происхождение признаков.
    
\end{abstract}

\textbf{Keywords:} TODO

\section{Introduction}\label{sec:intro}

TODO

\textbf{Contributions.} Our contributions can be summarized as follows:
\begin{itemize}
    \item We present...
    \item We demonstrate the validity of our theoretical results through empirical studies...
    \item We highlight the implications of our findings for...
\end{itemize}

\textbf{Outline.} The rest of the paper is organized as follows...

\section{Related Work}\label{sec:rw}

\textbf{Topic \#1.}
TODO

\textbf{Topic \#2.}
TODO

\section{Preliminaries}\label{sec:prelim}

\subsection{General notation}

In this section, we introduce the general notation used in the rest of the paper and the basic assumptions. 

\subsection{Assumptions} 

TODO

\section{Method}\label{sec:method}

\section{Experiments}\label{sec:exp}

To verify the theoretical estimates obtained, we conducted a detailed empirical study...

\section{Discussion}\label{sec:disc}

TODO

\section{Conclusion}\label{sec:concl}

TODO


%%%%%%%%%%%%%%%%%%%%%%%%%%%%%%%%%%%%%%%%%%%%%%%%%%%%%%%%%%%%

\bibliographystyle{unsrtnat}
\bibliography{references}

%%%%%%%%%%%%%%%%%%%%%%%%%%%%%%%%%%%%%%%%%%%%%%%%%%%%%%%%%%%%

\newpage
\appendix
\section{Appendix / supplemental material}\label{app}

\subsection{Additional experiments / Proofs of Theorems}\label{app:exp}

TODO

\end{document}
